\section{DEFINICIÓN DE OBJETIVOS}

\subsection{OBJETIVO GENERAL}

Diseñar un protocolo de votación electrónica basado en criptografía post-cuántica de retículos que mantenga la integridad y la confidencialidad del proceso de votación frente a amenazas de seguridad actuales y futuras.

\subsection{OBJETIVOS ESPECÍFICOS}
\begin{itemize}
	\item Fundamentar teóricamente el protocolo propuesto a partir de los estándares FIPS 203 y FIPS 204 y de la dificultad de los problemas Module-LWE y Module-SIS.
	\item Modelar formalmente el protocolo de votación electrónica mediante Tamarin Prover bajo el modelo de atacante Dolev-Yao.
	\item Verificar, mediante lemmas formales, las propiedades de confidencialidad, integridad y autenticidad del protocolo modelado.
	\item Determinar el rendimiento del protocolo propuesto en términos de latencia, tiempo de ejecución y consumo de memoria.
	\item Comparar el rendimiento y las propiedades de seguridad del protocolo propuesto con un protocolo equivalente basado en RSA/ECDSA.
\end{itemize}